Для практической верификации разработанной математической модели и~алгоритмов многозадачной координации разработана программное средство на языке программирования Go\footnote{\  исходный код и описание размещено на https://go.dev}. На программное средство получено Свидетельство о государственной регистрации программы для ЭВМ \cite{progbib1} (приложение~\ref{app:register_code}). Полный текст программного средства для самостоятельного проведения эксперимента или предложений для апробации доступен на https://github.com/dreddsa5dies/phd/Dissertation/code, а также в приложении~\ref{app:code}.

Выбор данного языка программирования обусловлен рядом его фундаментальных преимуществ, критически важных для построения надёжных, производительных и~масштабируемых систем управления роботизированными группами, как это детально описано в работах \cite{GoPractice, GoProfessional, GoLangDonovan, GoLangSumm}. Одновременно с этим, автор имеет большой опыт\footnote{\  программы с открытым исходным кодом размещены на https://github.com/dreddsa5dies, коммерческие продукты разрабатывались в компаниях Qtech и YADRO} в написании распределённых и~эффективных программных средств и~систем на этом языке программирования высокого уровня. Кроме озвученных фактов, имеется еще ряд безусловных преимуществ.

Так, Go изначально спроектирован для создания высоконагруженных сетевых сервисов и~распределенных систем. Имеет встроенную поддержку конкурентности через механизм горутин (goroutines) и каналов (channels) \cite{GoLangDonovan}, что позволяет просто и эффективно моделировать поведение множества автономных НТТС, действующих параллельно и взаимодействующих друг с другом. Это идеально соответствует задаче моделирования децентрализованной координации, где каждая машина принимает решения независимо, но с учетом состояния группы.

Во-вторых, язык обеспечивает исключительную производительность, близкую к C/C++, при этом сохраняя простоту синтаксиса и скорость разработки, характерную для интерпретируемых языков \cite{GoProfessional}. Это позволит реализовать вычислительно интенсивные симуляции с большим количеством итераций и машин за приемлемое время, что является ключевым требованием для проведения серийных экспериментов и получения статистически значимых результатов.

В-третьих, Go обладает превосходной поддержкой кроссплатформенной компиляции \cite{GoLangSumm}. Это означает, что разработанная симуляционная среда может быть собрана и запущена без изменений на различных операционных системах, что значительно упрощает ее распространение и использование, а~также потенциальную интеграцию с различными аппаратными платформами (см.~подраздел~\ref{subsec:ch1/sec2/sub4}).

Наконец, статическая типизация и строгая система компиляции Go \cite{GoPractice} способствуют созданию более надежного и предсказуемого кода. Это критически важно для научных исследований, где воспроизводимость результатов является основополагающим принципом. Компилятор выявляет множество потенциальных ошибок на этапе сборки, что снижает вероятность возникновения трудноуловимых багов в процессе выполнения симуляций.

Наиболее существенным недостатком использования этого языка программирования в части, касающейся применения в НТТС, является работа планировщика языка Go и использования механизма <<остановки>> для сборки неиспользуемых объектов в памяти ОС. Это значит, что его нельзя применять в~процессах, связанных с непосредственной опасностью для людей и окружающего мира. Для этого стоит применять только ОС реального времени и языки программирования без временных пауз. Но это не относится теме исследования и применение Go целесообразно для проведения эксперимента.

Общая архитектура реализации определена принципами модульности и~расширяемости. Система состоит из нескольких компонентов представленных на рисунке \ref{fig:program_modules}.

\begin{figure}[ht]
	\centering{
		\includegraphics[scale=0.6]{drawio/programm.drawio.png}
	}
	\caption{Модульная структура программной реализации метода автоматизированной многозадачной координации.}
	\label{fig:program_modules}
\end{figure}

Модуль стратегии (прил.~\ref{app:code8}) координации реализуют интерфейс Strategy (лист.~\ref{lst:strategy}), что позволяет легко добавлять новые алгоритмы распределения задач и сравнивать их в единой среде. Каждая стратегия должна удовлетворять требованиям этого интерфейса, т.\ е. запускаться, принимая набор машин и задач, и~возвращать метрики для последующего анализа. Метод String() предназначен для возврата наименования стратегии.

\begin{ListingEnv}[!ht]
	\captiondelim{ }
	\caption{Интерфейс стратегий по выбору задач}\label{lst:strategy}
	\begin{lstlisting}[language=Go]
		type Strategy interface {
			Run(machines []models.Machine, tasks []models.Task) StepMetrics
			String() string
		}
	\end{lstlisting}
\end{ListingEnv}%

В моделях данных определяются структуры данных для представления машин (прил.~\ref{app:code4}) и задач (прил.~\ref{app:code5}), включая их характеристики. Пример описания модели машин с распределением по типам представлен в описаниях \ref{lst:machine}, \ref{lst:type_machine}. Описаны такие типы, как экскаватор с набором оборудования: ковш, грейфер и~рыхлитель (типы 1, 3, 9), бульдозер с отвалом и рыхлителем (типы 2, 9), погрузчик с ковшом, захватом и вилами (типы 1, 7, 14) и другие.

\begin{ListingEnv}[!ht]
	\captiondelim{ }
	\caption{Модель данных машин}\label{lst:machine}
	\begin{lstlisting}[language=Go]
		// Machine описывает машину
		type Machine struct {
			// уникальный идентификатор
			ID string `json:"id" yaml:"id"`
			// тип оборудования
			Equipment []TypeEquipment `json:"equipment" yaml:"equipment"`
			// запас энергии
			Energy float64 `json:"energy" yaml:"energy"`
			// выполняемая задача
			AssignedTask *Task
		}
	\end{lstlisting}
\end{ListingEnv}%

\begin{ListingEnv}[!ht]
	\captiondelim{ }
	\caption{Описание типов оборудования}\label{lst:type_machine}
	\begin{lstlisting}[language=Go]
		const (
			Bucket     TypeEquipment = iota + 1 // ковш
			Dump                                // отвал
			Grapple                             // грейфер
			Brushes                             // щетки
			Cutter                              // фреза
			Plow                                // плуг
			Catch                               // захват
			Tamp                                // трамбовочное устройство
			Ripper                              // рыхлитель
			Scraper                             // скребок
			Conveyor                            // конвейер ленточный
			Knife                               // нож-отвал автогрейдера
			VCompactor                          // виброплита
			Pitchfork                           // вилы (для паллет)
			Thrower                             // метатель
		)
	\end{lstlisting}
\end{ListingEnv}%

Симуляционное ядро (прил.~\ref{app:code6}) отвечает за инициализацию эксперимента, последовательный запуск выбранной стратегии на заданном количестве шагов и~сбор метрик производительности.

Утилиты и конфигурация (прил.~\ref{app:code2} и~\ref{app:code3}) обеспечивают загрузку параметров эксперимента из внешнего файла типа YAML\footnote{\  формат описания данных, близкий к языкам разметки с поддержкой комментариев} и вспомогательные функции.

\nomenclature{YAML}{сокращенно от англ. yet another markup language -- еще один язык разметки\nomrefpage}

Аналитический модуль (прил.~\ref{app:code14}) позволяет генерировать сводные отчеты в формате CSV\footnote{\  простой табличный формат описания данных} и формирует результаты в виде графиков для наглядного сравнения эффективности различных стратегий.

\nomenclature{CSV}{сокращенно от англ. Comma-Separated Values -- значения, разделённые запятыми\nomrefpage}

Таким образом, реализация \cite{progbib1} на языке Go предоставляет надежную, производительную и гибкую платформу для проведения вычислительных экспериментов, что позволяет всесторонне оценить эффективность предложенного метода автоматизированной многозадачной координации по ключевым метрикам: времени выполнения задач, энергопотреблению и устойчивости к~сбоям связи, а также провести его сравнительный анализ с существующими аналогами.